\chapter{La valutazione e l'allenamento del metabolismo anaerobico}
\label{cap:metabolismoanaerobico}

Il modulo nasce da ciò che la match analysis (cap.~\ref{cap:matchanalysis}) ha misurato: stabilita
la richiesta fisiologica e cinematica della partita, si ricavano le \textbf{capacità fisiche
specifiche} del calcio a 5. Sono due, e sono quelle più rilevanti ai fini della prestazione:

\begin{itemize}
  \item l'\textbf{agility};
  \item la \textbf{repeated sprint ability} (RSA).
\end{itemize}

Si tratta prima l'\textbf{agility}, in entrambi i suoi versanti --- quello \textbf{atletico} e
quello \textbf{cognitivo} --- poi la \textbf{repeated sprint ability}.

% ---------------------------------------------------------------------------
\section{Che cos'è l'agility}

\rubrica{L'evoluzione della definizione}
La scienza dell'allenamento ha cambiato il concetto proprio a partire dalla sua definizione.
\begin{enumerate}
  \item \textbf{prima del 2002}: le definizioni guardano \textbf{esclusivamente} alla capacità di
    \textbf{cambiare direzione} in modo rapido;
  \item \textbf{Young e coll.} (2002): l'agilità è costituita da \textbf{due fattori
    fondamentali}, la \textbf{velocità nel cambio di direzione} e l'\textbf{aspetto cognitivo};
  \item \textbf{Sheppard e coll.} (2006): ``rapido movimento del corpo con cambio di velocità e
    direzione \textbf{in risposta a uno stimolo}''.
\end{enumerate}

Il passaggio decisivo è l'ingresso dell'\textbf{aspetto cognitivo} dagli anni Duemila. Sheppard in
particolare, parlando di risposta a uno stimolo, vi inserisce la \textbf{capacità di reazione}, che
è una delle \textbf{capacità coordinative}.

\rubrica{Le due componenti}
\begin{itemize}
  \item \textbf{performance atletica}: la capacità di cambiare direzione;
  \item \textbf{performance cognitiva}: la capacità di \textbf{percepire} --- di catturare
    informazioni dall'ambiente --- e di \textbf{decidere} la soluzione al problema che si presenta
    sul momento.
\end{itemize}

\rubrica{Perché è una capacità discriminante}
L'agility separa l'atleta di alto livello da quello di livello inferiore
(cap.~\ref{cap:profilogiocatore}), e lo fa quasi \textbf{a prescindere} dal dato cinematico di
gara: un giocatore capace di cambiare direzione, decelerare, frenare, fermarsi e ripartire in un
altro senso è efficace in partita per questo, prima ancora che per quanto corre.

% ---------------------------------------------------------------------------
\section{La performance atletica}

Schema di \textbf{Young} (2002), modificato in \textbf{Serpell e coll.} (2011). L'agility vi si
articola in due rami; questa lezione sviluppa il secondo.

\begin{enumerate}
  \item \textit{perceptual and decision making} --- il ramo cognitivo:
    \begin{itemize}
      \item \textit{visual scanning}: esplorazione visiva;
      \item \textit{knowledge of situations}: conoscenza delle situazioni di gioco;
      \item \textit{pattern recognition}: riconoscimento degli schemi;
      \item \textit{anticipation}: anticipazione.
    \end{itemize}
  \item \textit{change of direction speed} --- il ramo atletico, che dipende da quattro fattori:
    \begin{itemize}
      \item \textbf{antropometria}: le proporzioni del corpo;
      \item \textbf{velocità di sprint in linea} (\textit{straight sprinting speed}), che vi entra
        solo in parte;
      \item \textbf{tecnica} di esecuzione del cambio di direzione:
        \begin{itemize}
          \item adattamento della falcata per \textbf{accelerare e decelerare};
          \item inclinazione del corpo e postura (\textit{body lean and posture});
          \item \textbf{appoggio podalico} (\textit{foot placement});
        \end{itemize}
      \item \textbf{qualità muscolari degli arti inferiori}:
        \begin{itemize}
          \item \textbf{forza reattiva};
          \item forza e \textbf{potenza concentrica};
          \item \textbf{squilibrio fra arto destro e arto sinistro}.
        \end{itemize}
    \end{itemize}
\end{enumerate}

L'\textbf{appoggio podalico} merita attenzione a sé: è una caratteristica di gioco fondamentale nel
futsal, perché diventa \textbf{complementare al dribbling}.

% ---------------------------------------------------------------------------
\section{Che cosa correla davvero con il cambio di direzione}

Lo schema elenca i fattori da cui il cambio di direzione dipende, ma la letteratura che li ha messi
alla prova dà un esito \textbf{sorprendentemente magro}: le qualità fisiche che ci si aspetterebbe
decisive spiegano poco.

\rubrica{Forza massima degli arti inferiori}
\textbf{Nessuna} correlazione statisticamente significativa con il cambio di direzione:
\begin{itemize}
  \item \textbf{Markovic} (2007): squat con 1 RM eseguita al \textit{multipower};
  \item \textbf{Peterson e coll.} (2006): squat con 1 RM eseguita con bilanciere, cambio di
    direzione valutato con il \textbf{T-test}.
\end{itemize}

\rubrica{Potenza degli arti inferiori}
Correlazione \textbf{moderata} ($r \approx 0,4$, valore medio) fra altezza di salto misurata con il
salto verticale --- il \textit{counter movement jump} --- e cambio di direzione
(\textbf{Brughelli e coll.}, 2008).

\rubrica{Velocità di sprint in linea}
Correlazione \textbf{bassa} fra velocità lineare e capacità di eseguire cambi di direzione
(\textbf{Sheppard e Young}, 2006; \textbf{Young e coll.}, 1996).

\subsection{Accelerazione, velocità massima e agility sono qualità distinte}

\textbf{Little e Williams} (2005), su \textbf{106 calciatori professionisti} valutati su sprint di
10\,m (accelerazione), sprint lanciato di 20\,m (velocità massima) e prova di agility a zigzag.

\begin{itemize}
  \item le tre prestazioni risultano \textbf{tutte} correlate in modo significativo
    ($p < 0,0005$);
  \item ma i \textbf{coefficienti di determinazione} ($r^2$) sono bassissimi:
    \begin{itemize}
      \item accelerazione e velocità massima: \textbf{39\,\%};
      \item accelerazione e agility: \textbf{12\,\%};
      \item velocità massima e agility: \textbf{21\,\%};
    \end{itemize}
  \item se ne conclude che accelerazione, velocità massima e agility sono \textbf{qualità
    specifiche} e \textbf{relativamente indipendenti} fra loro.
\end{itemize}

Non è detto, insomma, che chi corre forte uno sprint in linea sui 30--40\,m ottenga un buon
risultato al T-test, e viceversa.

\subsection{E anche i loro allenamenti non si trasferiscono}

\textbf{Young, McDowell e Scarlett} (2001) hanno verificato se l'allenamento della velocità in
linea si trasferisca all'agility e viceversa. \textbf{36 soggetti}, valutati su uno sprint in linea
di 30\,m e su \textbf{sei prove di agility} con 2--5 cambi di direzione ad angoli diversi; due
sedute a settimana per \textbf{sei settimane}, con sprint in linea di 20--40\,m oppure sprint di
20--40\,m con 3--5 cambi di direzione di 100\textdegree.

\begin{itemize}
  \item l'allenamento della \textbf{velocità} migliora in modo significativo lo sprint in linea, ma
    dà guadagni \textbf{limitati} nelle prove di agility --- e \textbf{più il compito di agility è
    complesso, minore è il trasferimento};
  \item l'allenamento dell'\textbf{agility} migliora in modo significativo le prove di cambio di
    direzione, ma \textbf{non} lo sprint in linea.
\end{itemize}

\rubrica{La conseguenza metodologica}
Le due conclusioni convergono: se si vogliono indagare accelerazione, velocità massima e agility
servono \textbf{test specifici} per ciascuna, e per allenarle servono \textbf{mezzi distinti}. Non
si può dedurre l'una dall'altra né allenarle insieme.

% ---------------------------------------------------------------------------
\section{I test di valutazione}

Sono numerosi, e quelli che seguono indagano l'agility nel suo versante \textbf{fisico}, cioè la
sola capacità di cambiare direzione. In tutti la misura è il \textbf{tempo di percorrenza}: con le
\textbf{fotocellule} si ottiene una precisione scientifica, con il solo \textbf{cronometro} un dato
comunque utilizzabile.

\subsection{Illinois Agility Test}

Riferimento: \textbf{Miller M. e coll.} (2006). È il più noto.

\rubrica{Il percorso}
Rettangolo di \textbf{10\,m} di lunghezza per \textbf{5\,m} di larghezza, con \textbf{quattro coni}
allineati al centro e distanziati \textbf{3,3\,m} l'uno dall'altro; partenza e arrivo ai due angoli
dello stesso lato corto.

\rubrica{L'esecuzione}
\begin{enumerate}
  \item partenza \textbf{da terra}, con sprint alla massima velocità;
  \item primo cambio di direzione a circa \textbf{180\textdegree};
  \item corsa a \textbf{zigzag} fra i coni centrali, molto intensa;
  \item ultimo cambio di direzione a \textbf{180\textdegree} e arrivo.
\end{enumerate}

\subsection{505 test}

Riferimento: \textbf{Draper J.A.} e \textbf{Lancaster M.G.} (1985), quindi ben antecedente al
precedente.

\rubrica{Materiale}
Fotocellule, un cronometro, \textbf{sei coni} e un metro per misurare le distanze: è
\textbf{estremamente semplice da costruire} e molto pratico in campo.

\rubrica{Il percorso}
Dalla partenza si corre in linea per \textbf{10\,m} fino ai paletti, poi altri \textbf{5\,m} fino
al punto di \textbf{svolta}, dove si inverte il senso di marcia e si torna indietro. Valuta quindi
il \textbf{cambio di direzione a 180\textdegree}.

\rubrica{L'attenzione all'esecuzione}
Il test \textbf{dura pochissimo}, e per questo controllarne l'esecuzione è determinante: un
appoggio impreciso per il cambio di direzione --- eseguito \textbf{in anticipo} o
\textbf{eccessivamente in ritardo} --- produce un risultato non corretto.

\subsection{T-test}

Riferimento: \textbf{Pauole K. e coll.} (2000). Noto quanto l'Illinois.

\rubrica{Il percorso}
A forma di T: dalla linea di partenza/arrivo si trova il cono \textbf{A} a \textbf{5\,m}; i coni
\textbf{B} e \textbf{C} stanno ai lati di A, e la slide indica in \textbf{5\,m} la distanza
\textbf{da B a C}. Lo schema \textbf{quotato} ripreso nella parte del corso dedicata al calcio a 11
(cap.~\ref{cap:testvalutazionecalcio}) chiarisce che i \textbf{5\,m} vanno intesi \textbf{per
ciascun braccio} --- 4,57\,m, cioè 5 yard --- mentre per l'\textbf{asse} dalla partenza al cono
centrale dà \textbf{9,14\,m} (10 yard), quota che non coincide con i 5\,m indicati qui.

\rubrica{L'esecuzione}
\begin{enumerate}
  \item \textbf{sprint} fino al cono A, che va toccato;
  \item \textbf{corsa laterale} (\textit{slide}) verso destra fino al cono B, toccandolo;
  \item corsa laterale in senso opposto fino al cono C, toccandolo;
  \item corsa laterale di ritorno fino ad A, toccandolo;
  \item \textbf{corsa all'indietro} fino alla linea di partenza/arrivo.
\end{enumerate}

\rubrica{Che cosa lo rende interessante}
Alterna \textbf{modalità di corsa diverse} nello stesso test --- frontale, laterale, all'indietro
--- e richiede appoggi \textbf{molto forti} nei momenti di inversione del movimento.

\rubrica{I valori di riferimento}
Sono gli unici forniti nel corso per un test di agility, e arrivano dalla parte dedicata al calcio a
11 (cap.~\ref{cap:testvalutazionecalcio}): la scala di giudizio va da \textbf{eccellente}
($<$ 9,5$''$) a \textbf{scarso} ($>$ 11,6$''$), mentre i valori attesi per ruolo sono
$8{,}06 \pm 0{,}27$\,s per i \textbf{difensori}, $8{,}35 \pm 0{,}26$\,s per i
\textbf{centrocampisti} e $8{,}38 \pm 0{,}28$\,s per gli \textbf{attaccanti}.

\subsection{15-m Agility Run}

\rubrica{Il percorso}
Corsa su \textbf{15\,m} a \textbf{partenza lanciata}: il cronometro parte dal segno dello
\textbf{0\,m}, preceduto da \textbf{3\,m} di rincorsa.
\begin{itemize}
  \item dallo 0\,m, \textbf{3\,m} di corsa libera;
  \item \textbf{3\,m} di \textbf{slalom} fra tre sagome;
  \item \textbf{2\,m} fino a un \textbf{ostacolo} (\textit{hurdle}) da superare;
  \item \textbf{7\,m} finali fino al traguardo dei 15\,m.
\end{itemize}

\subsection{Change of Direction Speed Test}

Riferimento: \textbf{Sheppard J.M. e coll.} (2006). Dalla partenza il soggetto raggiunge, su un
tracciato complessivo di \textbf{10\,m}, uno di due cancelletti laterali. Il cambio di direzione è
\textbf{volontario}: il lato è deciso in anticipo.

\subsection{I test con componente reattiva}

Qui il cambio di direzione non è più deciso dal soggetto ma imposto da uno \textbf{stimolo
esterno}. È il confronto fra le due forme --- cambio di direzione \textbf{volontario} contro cambio
di direzione \textbf{in funzione di uno stimolo dato da un operatore} --- che comincia a portare la
valutazione sul versante \textbf{cognitivo}.

\rubrica{Reactivity Agility Sprint Speed (Oliver J.L. e Meyers R.W., 2009)}
\begin{itemize}
  \item alla partenza, cellula fotoelettrica e cellula riflettente poste a \textbf{1\,m};
  \item \textbf{5\,m} di corsa fino a un \textbf{cancelletto temporizzato centrale}, fiancheggiato
    da barriere di gommapiuma (70\,cm di altezza $\times$ 90 di lunghezza $\times$ 30 di
    larghezza);
  \item dal cancelletto si aprono \textbf{tre uscite}: \textit{left}, \textit{straight},
    \textit{right}. Quella diritta è a \textbf{5\,m}; le due laterali sono anch'esse a 5\,m, come
    ipotenusa di un percorso di 4\,m in avanti e 3\,m di lato.
\end{itemize}

\rubrica{Reactive Agility Test --- RAT (Sheppard J.M. e coll., 2006)}
\begin{itemize}
  \item il partecipante parte da una linea e corre verso il centro;
  \item ai due lati stanno altrettanti \textbf{cancelletti temporizzati}, profondi \textbf{2\,m};
  \item l'\textbf{operatore} parte da un \textbf{tappeto temporizzato} posto al centro e, con il
    proprio movimento, fornisce lo \textbf{stimolo}: il partecipante deve riconoscerlo e scegliere
    di conseguenza il lato verso cui uscire.
\end{itemize}

\subsection{Il limite di questi test}

Indagano \textbf{esclusivamente} la parte del cambio di direzione, e non la componente cognitiva.

Un'integrazione a costo quasi nullo consiste nel \textbf{riprendere il test con una telecamera}:
il filmato permette di valutare anche le componenti \textbf{coordinative} --- il modo in cui il
giocatore affronta il cambio di direzione --- e di far emergere l'\textbf{arto dominante}, cioè
quale dei due spinge di più. Sono indicazioni direttamente utili al preparatore, perché dicono
\textbf{su quale arto} intervenire con il lavoro di forza.

% ---------------------------------------------------------------------------
\section{La performance cognitiva}

È il ramo \textit{perceptual and decision making} dello schema di Young: esplorazione visiva,
conoscenza delle situazioni, riconoscimento degli schemi, anticipazione.

\rubrica{Perché il calcio lo richiede}
Come \textbf{sport di situazione}, il calcio prevede abilità tecniche di natura \textbf{``open''}
piuttosto che \textbf{``closed''}: la prestazione si configura mediante una serie di operazioni
\textbf{sia mentali sia motorie}, e l'\textbf{anticipazione} e la \textbf{percezione dei dettagli
dell'ambiente} sono fondamentali tanto per i processi decisionali ed esecutivi quanto per quelli
interpretativi (\textbf{D'Ottavio}, 2011).

\rubrica{Il terzo stadio dell'evoluzione}
Al percorso già visto --- cambio di direzione, poi risposta a uno stimolo --- si aggiunge ora
l'\textbf{anticipazione}: non più solo reagire a uno stimolo \textbf{una volta che si è
presentato}, ma \textbf{prevedere ciò che sta per avvenire} e decidere di conseguenza.

\subsection{L'anticipazione}

\textbf{Anticipazione}: il piano di organizzazione \textbf{mentale e motoria} elaborato in termini
di \textbf{probabilità}; la capacità di vedere e conoscere \textbf{in anticipo} il senso
dell'azione o di un comportamento tecnico (\textbf{D'Ottavio}, 2011).

\rubrica{Su che cosa si fonda}
Sull'analisi dei \textbf{``segni''} e delle \textbf{``specificità''}:
\begin{itemize}
  \item del \textbf{gioco} e delle sue situazioni;
  \item degli \textbf{esercizi};
  \item dei \textbf{compagni di squadra} --- che cosa sta per fare il compagno, sul piano tattico
    (senza palla) e tecnico (con palla);
  \item degli \textbf{avversari};
  \item della \textbf{posizione della palla}, per esempio per prevederne la traiettoria e
    intercettarla.
\end{itemize}

\rubrica{Il peso dell'esperienza}
Buona parte della capacità di anticipazione dipende dall'\textbf{esperienza di gioco}. Dal
giocatore d'élite ci si attende perciò una catena completa e rapida: \textbf{ricevere}
informazioni dall'ambiente $\rightarrow$ \textbf{selezionare} le più importanti $\rightarrow$
\textbf{decidere} su base probabilistica $\rightarrow$ \textbf{eseguire} il gesto.

\subsection{Il fattore tempo}

Nell'anticipazione il tempo è la variabile decisiva: non conta soltanto \textbf{saper} eseguire un
gesto --- un passaggio, una trasmissione --- ma \textbf{quanto velocemente} lo si esegue.

\rubrica{Il conflitto fra velocità e precisione}
Le due qualità si contendono lo stesso margine: \textbf{più il gesto va eseguito in fretta, più si
rischia di essere imprecisi}. Il \textbf{talento sportivo} si riconosce proprio qui --- è chi
riesce a eseguire il gesto tecnico in modo \textbf{estremamente rapido} mantenendolo nella forma
\textbf{più precisa possibile}.

\rubrica{In allenamento}
\begin{itemize}
  \item \textbf{avere meno tempo} per analizzare la situazione;
  \item quindi \textbf{velocità mentale} richiesta \textbf{sia prima sia durante} l'esecuzione ---
    non solo per via verbale, ma per come l'esercizio stesso è costruito.
\end{itemize}

\rubrica{In gara}
La stessa compressione del tempo è determinata dalla \textbf{pressione} e dalla \textbf{presenza
dell'avversario}.

\subsection{Come si allena l'anticipazione}

Si organizzano esercizi che intervengono sui due momenti del processo:

\begin{enumerate}
  \item esercizi che \textbf{limitano o influenzano la percezione e l'attenzione};
  \item esercizi che \textbf{limitano o influenzano l'elaborazione}, cioè il momento della
    decisione (\textit{decision making}).
\end{enumerate}

In concreto, esercizi che:
\begin{itemize}
  \item stimolano la \textbf{velocità di percezione};
  \item stimolano la \textbf{velocità di decisione};
  \item stimolano la \textbf{riorganizzazione rapida} in base a eventi non previsti --- per
    esempio situazioni di gioco in cui il senso dell'azione cambia a un segnale.
\end{itemize}

\subsection{I tre livelli di interferenza cognitiva}

Classificazione di \textbf{D'Ottavio} (2011), che ordina le prove secondo quanto carico cognitivo
impongono.

\begin{enumerate}
  \item \textbf{livello 0} --- \textbf{COD}, nessuna interferenza cognitiva: programmazione di
    un'azione motoria \textbf{in assenza di stimoli esterni};
  \item \textbf{livello 1} --- \textbf{agility} con interferenza cognitiva di primo livello
    (\textbf{S-R}): programmazione di un'azione motoria \textbf{in risposta a uno stimolo esterno}
    non preceduto da preavviso (stimolo $\rightarrow$ risposta);
  \item \textbf{livello 2} --- \textbf{agility} con interferenza cognitiva di secondo livello
    (\textbf{percezione e analisi della situazione}): programmazione di un'azione motoria
    \textbf{in seguito all'elaborazione di un'informazione esterna} (situazione $\rightarrow$
    anticipazione).
\end{enumerate}

La differenza fra livello 1 e livello 2 è decisiva: nel primo il giocatore \textbf{reagisce} a uno
stimolo, nel secondo \textbf{legge} una situazione in corso e ne \textbf{prevede} l'esito.

\subsection{Lo studio pilota sui tre livelli}

Da questi presupposti è stata costruita una sequenza di valutazione che porta l'anticipazione
dentro la misura dell'agility, provata sul campo con la \textbf{nazionale di calcio a 5}.

\rubrica{Il dispositivo}
Quattro fotocellule: \textbf{FC1} alla partenza, \textbf{FC2} nel punto in cui va eseguito il
cambio di direzione, \textbf{FC3} e \textbf{FC4} sulle due uscite, a destra e a sinistra. Il
giocatore parte da fermo, sprinta alla massima velocità e cambia direzione all'altezza di FC2.

\rubrica{Livello 0 --- cambio di direzione puro}
Partenza \textbf{volontaria}, senza segnale di via: è il passaggio del giocatore attraverso la
porta a far partire il cronometro. Nessuno stimolo esterno; il lato del cambio di direzione è
stabilito in anticipo, con due prove, una a destra e una a sinistra.

\rubrica{Livello 1 --- reattività}
All'altezza di FC2 si accende un \textbf{LED}, collegato direttamente alla fotocellula:
\begin{itemize}
  \item \textbf{verde} $\rightarrow$ cambio di direzione a \textbf{destra};
  \item \textbf{rosso} $\rightarrow$ cambio di direzione a \textbf{sinistra}.
\end{itemize}
Lo stimolo arriva \textbf{nell'istante} del passaggio: il giocatore non può prepararsi.

\rubrica{Livello 2 --- ``read and react''}
Mentre il giocatore corre, \textbf{due operatori si passano la palla} in forma standardizzata:
a \textbf{5\,m} di distanza fra loro, cioè 2,5\,m per parte rispetto all'asse, con ritmo il più
regolare possibile e \textbf{senza finte}. Al passaggio su FC2 il giocatore va \textbf{dalla parte
dove si trova la palla}. La differenza rispetto al livello 1 è che qui, già durante il primo tratto
di sprint, il giocatore \textbf{può prevedere} dove sarà la palla al momento del suo passaggio, e
decidere \textbf{in anticipo}.

\subsection{I risultati}

\begin{table}[htbp]
  \centering
  \small
  \renewcommand{\arraystretch}{1.3}
  \begin{tabularx}{\textwidth}{|X|c|c|c|c|c|c|}
    \hline
    & \multicolumn{2}{c|}{\textbf{COD}} & \multicolumn{2}{c|}{\textbf{COD reactivity}} &
      \multicolumn{2}{c|}{\textbf{Read \& react}} \\
    \cline{2-7}
    & dx & sx & dx & sx & dx & sx \\
    \hline
    N.\ prove & 3 & 3 & 2 & 2 & 5 & 2 \\
    \hline
    Migliore (s) & 2,595 & 2,636 & 3,190 & 3,176 & 2,572 & 2,642 \\
    \hline
    Media (s) & 2,659 & 2,738 & 3,223 & 3,634 & 2,843 & 2,765 \\
    \hline
    DS & 0,06 & 0,09 & 0,05 & 0,65 & 0,30 & 0,17 \\
    \hline
  \end{tabularx}
\end{table}

\rubrica{Il costo della reazione}
Passando dal livello 0 al livello 1 la prestazione \textbf{peggiora nettamente}:
\textbf{$-$22,9\,\%} a destra e \textbf{$-$20,5\,\%} a sinistra. Quello scarto \textbf{è} il
\textbf{tempo di reazione}: è ciò che il giocatore paga per dover attendere lo stimolo.

\rubrica{Che cosa fa l'anticipazione}
Passando dal livello 0 al livello 2 lo scarto quasi scompare: \textbf{$+$0,9\,\%} a destra e
\textbf{$-$0,2\,\%} a sinistra, cioè una differenza \textbf{inferiore all'1\,\%} fra l'esecuzione
\textbf{puramente atletica} e quella che richiede la capacità di anticipazione. Nel cambio di
direzione a destra il tempo con anticipazione è addirittura \textbf{migliore} della miglior prova
senza interferenze cognitive.

\rubrica{La conclusione}
L'\textbf{anticipazione azzera}, o comunque riduce in misura importante, il \textbf{tempo di
reazione}. È la ragione per cui una valutazione dell'agility fermata al livello 0 --- quella dei
test del paragrafo precedente --- misura solo una parte del fenomeno.

% ---------------------------------------------------------------------------
\section{La repeated sprint ability}

È la seconda delle due capacità specifiche del futsal.

\subsection{Le definizioni}

\begin{itemize}
  \item \textbf{Bishop}: l'abilità di produrre la \textbf{miglior prestazione media} su una serie
    di sprint ($\leq$ 10\,s), separati da brevi periodi di recupero ($\leq$ 60\,s);
  \item \textbf{Thébault}: l'abilità di \textbf{ripetere} brevi periodi di sprint con un breve
    recupero fra essi;
  \item \textbf{Bishop}, in una formulazione successiva: l'abilità di fornire prestazioni di
    sprint con un \textbf{minor decremento} della prestazione massima --- l'accento si sposta
    sulla \textbf{produzione meccanica} della prestazione;
  \item \textbf{D'Ottavio}: l'abilità di reiterare sprint con ridotto decremento della
    prestazione, cioè la capacità di \textbf{sprintare, accelerare e svolgere movimenti brevi ad
    alta intensità, quindi recuperare e sprintare ancora}.
\end{itemize}

L'ultima è la più completa, perché non descrive solo il decremento ma l'intero \textbf{modello
dell'alta intensità} negli sport di squadra: azioni brevi e intense in spazi ristretti, recupero
il più breve possibile, e di nuovo azione.

\subsection{Perché è specifica del futsal}

La match analysis (cap.~\ref{cap:matchanalysis}) l'ha già misurata: nel calcio a 5 i giocatori
ripetono due, tre, quattro sprint quasi senza soluzione di continuità, con recuperi non solo
inferiori ai 60\,s ma spesso di soli \textbf{15\,s} --- la sequenza in assoluto più frequente è
proprio \textbf{due sprint separati da 15 secondi}, e si registrano anche sequenze di quattro
sprint con lo stesso intervallo.

Essere performanti in quella serie e insieme recuperare nel minor tempo possibile è quindi una
capacità \textbf{fondamentale} per questo sport, e come l'agility \textbf{discrimina} il giocatore
d'élite da quello che non lo è. Va inoltre letta dentro uno sport a \textbf{\textit{open skill}},
dove all'esecuzione si somma la componente cognitiva già vista --- reazione e anticipazione.

\subsection{L'età: quando conviene allenarla}

\rubrica{Mujika e coll. (2009)}
\begin{itemize}
  \item \textbf{134 giovani calciatori} distribuiti per fascia d'età, dagli U-11 agli U-18
    ($n$ fra 11 e 22 per categoria);
  \item test di RSA \textbf{6 $\times$ 30\,m} con recupero \textbf{attivo} di 30$''$;
  \item variabili: \textbf{tempo totale} (TT) e \textbf{percentuale di decremento} (Dec\%);
  \item la prestazione nel TT \textbf{migliora durante la maturazione}, in modo marcato dagli
    \textbf{U-11 agli U-15} ($p < 0,05$) e più lentamente in seguito ($p < 0,05$).
\end{itemize}

\rubrica{Wierike e coll. (2013)}
\begin{itemize}
  \item \textbf{48 giocatori di basket d'élite} fra i 14 e i 19 anni, testati in \textbf{6
    occasioni} nelle stagioni 2008-09 e 2009-10;
  \item la RSA \textbf{migliora con l'età}, soprattutto fra i \textbf{14 e i 17 anni}
    ($p < 0,05$), raggiungendo un \textbf{plateau} fra i 17 e i 19.
\end{itemize}

\rubrica{La fase sensibile}
I due studi convergono: la finestra utile è quella delle categorie \textbf{U-11--U-15}, cioè fra i
\textbf{10 e i 14 anni} d'età, che il secondo studio estende fino ai 17. In quella fascia il
sistema è \textbf{pronto} a ricevere questo tipo di sollecitazione e vi si adatta in modo
particolarmente efficace.

\subsection{Il livello di qualificazione}

\begin{itemize}
  \item \textbf{Aziz e coll.} (2008), su professionisti, semiprofessionisti e amatori: la
    prestazione di RSA è \textbf{correlata al livello di competizione}, con differenze
    statisticamente significative;
  \item \textbf{Rampinini e coll.} (2009), professionisti contro amatori: differenze
    statisticamente significative fra i due livelli, con \textbf{indici di affaticamento} diversi;
  \item \textbf{Gabbett} (2010), su giocatrici d'élite di livello \textbf{nazionale} e di
    \textbf{club} ($n = 19$; 18,1 $\pm$ 2,9 anni): le prove di sprint ripetuti
    \textbf{discriminano} i due livelli.
\end{itemize}

\rubrica{Che cosa se ne ricava per il campo}
\begin{itemize}
  \item allenando giocatori di \textbf{alta qualificazione} si parte da una base già molto alta,
    perché la RSA è intrinseca al gioco ed è verosimilmente stata un \textbf{elemento di
    selezione} per l'accesso a quel livello;
  \item ai \textbf{livelli più bassi} va invece \textbf{strutturata e migliorata}, non soltanto
    mantenuta: servono strategie di allenamento apposite.
\end{itemize}

\subsection{Stretching e riscaldamento}

La domanda operativa è se lo \textbf{stretching statico} inserito nel riscaldamento comprometta la
prestazione successiva --- questione aperta soprattutto per le sedute con obiettivo di forza.

\rubrica{Sim e coll. (2009)}
13 giocatori di sport di squadra, protocollo di \textbf{6 $\times$ 20\,m} con 25$''$ di recupero,
preceduto in giorni diversi da tre riscaldamenti:
\begin{itemize}
  \item \textbf{D}: 1000\,m di jogging seguiti da attività dinamiche;
  \item \textbf{SD}: stretching statico seguito da attività dinamiche;
  \item \textbf{DS}: attività dinamiche seguite da stretching statico.
\end{itemize}
Variabili analizzate: primo sprint (FST), sprint migliore (BST) e tempo totale (TST).

\begin{itemize}
  \item i tempi sono risultati \textbf{più lenti} nella condizione \textbf{DS};
  \item ma \textbf{nessuna} differenza statisticamente significativa fra TST, BST e FST, né fra DS
    e SD;
  \item gli autori concludono ugualmente che la RSA \textbf{potrebbe} essere compromessa da un
    riscaldamento di tipo DS --- una tendenza che i dati suggeriscono \textbf{senza} che
    l'evidenza statistica la sostenga.
\end{itemize}

\rubrica{Wong e coll. (2011)}
Su calciatori di 16,8 $\pm$ 0,4 anni: \textbf{tre giorni consecutivi} di stretching statico
\textbf{non} hanno influenzato negativamente la RSA.

\rubrica{Turki-Belkhiria e coll. (2014)}
\textit{Eight weeks of dynamic stretching during warm-ups improves jump power but not repeated or
single sprint performance}: programmi con stretching statico e dinamico ripetuti per \textbf{otto
settimane} non producono effetti significativi né sulla RSA né sul singolo sprint.

\rubrica{Taylor e coll. (2013)}
11 calciatori \textit{subelite}, RSA su \textbf{6 $\times$ 40\,m} con 20$''$ di recupero, dopo tre
riscaldamenti:
\begin{enumerate}
  \item 5$'$ di jogging al 65\,\% della $FC_{max}$, seguiti da \textbf{attività specifica} con
    palla;
  \item 5$'$ di jogging al 65\,\% della $FC_{max}$, seguiti da \textbf{stretching statico} e poi
    da attività specifica;
  \item 5$'$ di jogging al 65\,\% della $FC_{max}$, seguiti da \textbf{stretching}.
\end{enumerate}

Nel grafico dei risultati le tre curve sono etichettate \textit{no stretching}, \textit{static
stretching} e \textit{dynamic stretching}.
\begin{itemize}
  \item in tutte e tre le condizioni i tempi \textbf{salgono progressivamente} dal primo al sesto
    sprint --- da circa 7,2 a circa 7,9\,s --- come atteso per l'affaticamento;
  \item la condizione con \textbf{stretching statico} è \textbf{costantemente la più lenta};
  \item le altre due sono \textbf{sostanzialmente sovrapponibili}.
\end{itemize}

Anche qui, dunque, lo stretching statico prima della prestazione sembra avere un'influenza
\textbf{non favorevole} sulla RSA.

\subsection{I test di valutazione}

Non esiste un test \textbf{univoco} con distanze standard: a modelli di prestazione diversi devono
corrispondere modi diversi di valutare la RSA. Fra un protocollo e l'altro cambiano la
\textbf{modalità} (corsa in linea, cicloergometro, navetta), la \textbf{distanza del singolo
sprint}, il \textbf{volume totale} e i \textbf{tempi di recupero}.

\rubrica{Repeated Sprint Ability --- modalità in linea}
\begin{itemize}
  \item 8 $\times$ 35\,m con 30$''$ di recupero (\textbf{Rushall e coll.}, 1991);
  \item 12 $\times$ 20\,m con 20$''$ di recupero (\textbf{Wadley e coll.}, 1998);
  \item 6 $\times$ 15$''$ su \textbf{cicloergometro} con 90$''$ di recupero
    (\textbf{McMahon e coll.}, 1998);
  \item 7 $\times$ 30\,m con 20$''$ di recupero (\textbf{Reilly e coll.}, 2000);
  \item 6 $\times$ 4$''$ con 25$''$ di recupero (\textbf{Spencer e coll.}, 2002).
\end{itemize}

\rubrica{Repeated Shuttle Sprint Ability --- modalità a navetta}
Introduce una componente neuromuscolare in più, il \textbf{cambio di direzione a
180\textdegree}.
\begin{itemize}
  \item 10 $\times$ 15\,m con 30$''$ di recupero (\textbf{Castagna e coll.}, 2005);
  \item 6 $\times$ 20\,m con 20$''$ di recupero (\textbf{Impellizzeri e coll.}, 2008);
  \item 6 $\times$ 40\,m con 20$''$ di recupero (\textbf{Rampinini e coll.}, 2009);
  \item 7 $\times$ 15\,m con recupero \textbf{1:5} (\textbf{Ruscello, D'Ottavio e coll.}, 2013).
\end{itemize}

\rubrica{Quale scegliere per il futsal}
È il \textbf{modello di prestazione} a guidare la scelta --- e a indicare, di conseguenza, anche i
mezzi di allenamento.
\begin{itemize}
  \item la \textbf{navetta} è più specifica della corsa in linea, perché impone almeno un cambio
    di direzione;
  \item i due protocolli su \textbf{15\,m} --- Castagna 2005 e Ruscello 2013 --- sono quelli più
    vicini alla \textbf{distanza media dello sprint} rilevata in partita, fra i 10 e i 15\,m
    (cap.~\ref{cap:matchanalysis});
  \item resta però un limite: il giocatore di futsal compie cambi di direzione a
    \textbf{180\textdegree} piuttosto \textbf{di rado}. Da qui la domanda se non convenga valutare
    la RSA attraverso \textbf{cambi di direzione} (COD) veri e propri.
\end{itemize}

\subsection{Il rapporto lavoro:recupero}

Studio di \textbf{Ruscello e coll.} (2013), \textit{Influence of the number of trials and the
exercise to rest ratio in repeated sprint ability, with changes of direction and orientation}.

\rubrica{Ipotesi}
Nell'allenamento della RSA il \textbf{numero di ripetizioni} e il \textbf{rapporto
lavoro:recupero} \textbf{non sono equivalenti} se applicati a modalità di sprint diverse. Le
domande di ricerca sono quindi quali siano il numero di ripetizioni e il rapporto ottimali per
ciascuna modalità.

\rubrica{Le quattro fasi}
\begin{enumerate}
  \item \textbf{ipotesi} e domande di ricerca;
  \item \textbf{raccolta dati} su 17 giocatori di calcio a 5, con protocollo a
    \textit{latin square} (abc; bca; cab) e test svolti in tre giorni diversi, tutti al rapporto
    \textbf{1:5};
  \item \textbf{analisi statistica e modello matematico}: ANOVA a misure ripetute, ANOVA
    fattoriale e analisi di regressione. Variabile dipendente l'\textbf{IF\%}, variabili
    indipendenti le tre modalità di sprint; il modello restituisce numero di ripetizioni e
    rapporto ottimali;
  \item \textbf{studio pilota} che applica in campo i rapporti ricavati --- linea 1:5, navetta
    1:3, COD 1:2 --- su tre giornate di test separate da 48 ore.
\end{enumerate}

\rubrica{Le tre modalità, a parità di distanza}
Il singolo sprint copre sempre \textbf{30\,m}, ripetuti \textbf{7 volte} per un totale di
\textbf{210\,m}; cambia solo come quei 30\,m vengono percorsi.
\begin{itemize}
  \item \textbf{linea}: 7 $\times$ 30\,m;
  \item \textbf{navetta}: 7 $\times$ (15 $+$ 15\,m);
  \item \textbf{COD}: 7 $\times$ (6 $\times$ 5\,m), a zigzag con angoli di 60\textdegree.
\end{itemize}

\rubrica{L'indice di fatica}
$$IF\% = \left(1 - \frac{[\text{personal best}]}{[\text{trial}]}\right) \times 100$$

\begin{table}[htbp]
  \centering
  \small
  \renewcommand{\arraystretch}{1.3}
  \begin{tabularx}{\textwidth}{|X|c|c|c|}
    \hline
    \textbf{IF\% per prova} & \textbf{Linea} & \textbf{Navetta} & \textbf{COD} \\
    \hline
    Prova 1 & 0,28 & 1,18 & 1,03 \\
    \hline
    Prova 2 & 2,30 & 1,70 & 1,38 \\
    \hline
    Prova 3 & 3,61 & 1,87 & 1,38 \\
    \hline
    Prova 4 & 4,88 & 2,87 & 1,73 \\
    \hline
    Prova 5 & 5,71 & 4,50 & 2,66 \\
    \hline
    Prova 6 & 5,91 & 6,85 & 5,00 \\
    \hline
    Prova 7 & 8,69 & 7,46 & 5,76 \\
    \hline
    \textbf{Media} & \textbf{4,48} & \textbf{3,78} & \textbf{2,71} \\
    \hline
    DS & 0,03 & 0,03 & 0,02 \\
    \hline
  \end{tabularx}
\end{table}

\rubrica{I punti di cut-off}
La ripetizione a partire dalla quale il decremento diventa significativo:
\begin{itemize}
  \item \textbf{linea}: \textbf{2/7} ($p = 0,014$) --- arriva \textbf{molto presto};
  \item \textbf{navetta}: 4/7 ($p = 0,004$);
  \item \textbf{COD}: 5/7 ($p = 0,020$).
\end{itemize}

A parità di rapporto 1:5, dunque, lo sforzo dello sprint \textbf{in linea} risulta il
\textbf{più impegnativo} sul piano del dispendio energetico, e quello con cambi di direzione il
meno.

\begin{table}[htbp]
  \centering
  \small
  \renewcommand{\arraystretch}{1.3}
  \begin{tabularx}{\textwidth}{|X|c|c|c|c|c|}
    \hline
    \textbf{Test} (7 $\times$ 30\,m, 1:5) & \textbf{IF\%} & \textbf{MTP} (s) &
      \textbf{MRT} (s) & \textbf{IIR} & \textbf{OTR} (s) \\
    \hline
    Linea & 4,48 & 4,53 & 22 & 0,204 & $\approx$22 (1:5) \\
    \hline
    Navetta & 3,78 & 5,90 & 29 & 0,130 & $\approx$14 ($\approx$1:2,5) \\
    \hline
    COD & 2,71 & 8,50 & 41 & 0,066 & $\approx$10 ($\approx$1:1) \\
    \hline
  \end{tabularx}
\end{table}

\begin{itemize}
  \item \textbf{MTP}, \textit{mean time of performance}: il tempo di esecuzione cresce molto ---
    gli stessi 30\,m costano 4,53\,s in linea e 8,50\,s a cambi di direzione;
  \item \textbf{MRT}, \textit{mean recovery time}: al rapporto 1:5 il recupero effettivo sale di
    conseguenza, da 22 a 41\,s;
  \item \textbf{IIR}, \textit{index of influence of recovery};
  \item \textbf{OTR}, \textit{optimal time of recovery}: il recupero che servirebbe per ottenere
    lo \textbf{stesso} modello di affaticamento della modalità in linea. Per la navetta non 29\,s
    ma \textbf{14}, per il COD appena \textbf{10}.
\end{itemize}

\rubrica{La correzione del modello}
I tempi ottimali grezzi vengono corretti con un fattore ricavato dal rapporto fra i tempi di
esecuzione delle modalità:
\begin{itemize}
  \item navetta: fattore \textbf{0,232}, tempo corretto $\approx$\textbf{18} $\pm$ 2\,s,
    cioè $\approx$\textbf{1:3};
  \item COD: fattore \textbf{0,467}, tempo corretto $\approx$\textbf{12} $\pm$ 2\,s,
    cioè $\approx$\textbf{1:2}.
\end{itemize}

\rubrica{Quante ripetizioni servono al rapporto 1:5}
Rovesciando la domanda --- se \textbf{tengo} il rapporto 1:5, quante ripetizioni servono per
raggiungere lo stesso affaticamento della modalità in linea? --- l'estrapolazione delle rette di
regressione fino alla decima prova indica che nelle modalità non lineari occorre superare le
\textbf{10 ripetizioni}.

\rubrica{Che cosa se ne ricava}
\begin{itemize}
  \item il \textbf{rapporto lavoro:recupero ottimale} dipende dalla modalità: \textbf{1:5} in
    linea, \textbf{1:3} a navetta, \textbf{1:2} con cambi di direzione;
  \item in alternativa, mantenendo il rapporto 1:5 nella modalità COD occorre eseguire
    \textbf{10--12 ripetizioni};
  \item per il calcio a 5 la modalità più vicina al modello di prestazione è quella a
    \textbf{cambi di direzione}, da somministrare però con il \textbf{recupero corretto} --- che a
    sua volta è quello più simile a ciò che accade in partita. Il rapporto \textbf{non} è un tempo
    fisso: il recupero discende dal tempo di attività.
\end{itemize}

% ---------------------------------------------------------------------------
\section{La fatica muscolare}

\textbf{Fatica muscolare}: declino della prestazione indotto dall'\textbf{esercizio} --- inteso sia
come esercizio allenante sia come prestazione di gara.

Chiude il modulo perché è il rovescio delle due capacità appena viste: dopo aver stabilito come si
valutano e si allenano, resta da capire come \textbf{contrastare} gli effetti che la prestazione
stessa produce sulla fisiologia del giocatore.

\rubrica{Fattori e sintomi}
La trattazione separa i due piani, e la distinzione è operativa:
\begin{itemize}
  \item i \textbf{sintomi} sono ciò che si riesce a \textbf{rilevare} --- con la ricerca di
    laboratorio, l'analisi dell'allenamento e la match analysis
    (cap.~\ref{cap:matchanalysis});
  \item i \textbf{fattori} sono il \textbf{perché} quei sintomi compaiono. Si interviene su questi
    per contrastare quelli.
\end{itemize}

\rubrica{Dipendenza dal compito}
I meccanismi della fatica sono \textbf{sport-specifici}. Negli sport di squadra la fatica nasce da
una prestazione \textbf{intermittente}, con momenti di alta intensità alternati a recuperi: nel
calcio a 5 con una frequenza cardiaca di gara intorno al \textbf{90\,\% della $FC_{max}$} e una
consistente produzione di lattato. La fatica di un centometrista è cosa del tutto diversa da quella
di un maratoneta.

\subsection{I fattori della fatica periferica}

Sono disposti in cerchio, senza gerarchia di importanza:
\begin{itemize}
  \item $\downarrow$ \textbf{glucosio};
  \item $\downarrow$ \textbf{glicogeno} --- la deplezione del glicogeno muscolare è documentata
    nel calcio;
  \item $\downarrow$ \textbf{ossigeno};
  \item $\downarrow$ \textbf{$H_2O$}, cioè disidratazione;
  \item $\uparrow$ \textbf{$H^+$};
  \item $\uparrow$ \textbf{$K^+$}, intra ed extracellulare;
  \item $\uparrow$ \textbf{temperatura} ambientale.
\end{itemize}

\subsection{I sintomi}

\rubrica{Riduzione della performance}
\begin{itemize}
  \item $\downarrow$ distanza percorsa --- assoluta o relativa al tempo;
  \item $\downarrow$ tempo trascorso ad alta velocità;
  \item $\uparrow$ tempo trascorso stando fermo;
  \item $\downarrow$ tempo speso ad alta accelerazione e decelerazione;
  \item $\downarrow$ contrasti nel secondo tempo.
\end{itemize}
Si registrano soprattutto nelle \textbf{parti finali} dei due tempi. Va però distinto il dato
cinematico da eventuali \textbf{scelte tattiche} che riducono corsa e alta intensità per altre
ragioni.

\rubrica{Alterata o ridotta esecuzione tecnica}
\begin{itemize}
  \item calcio: $\downarrow$ tiro in porta, $\downarrow$ passaggio, $\downarrow$ precisione di
    passaggio e tiro;
  \item rugby: $\downarrow$ intensità dei contatti, $\downarrow$ precisione del passaggio.
\end{itemize}

\rubrica{Misurazioni durante i test di valutazione}
Confrontando con un test di riferimento svolto a inizio stagione si rilevano:
\begin{itemize}
  \item $\uparrow$ tempi negli sprint ripetuti e nei test di agility;
  \item $\downarrow$ picco di forza isocinetica;
  \item $\downarrow$ potenza degli arti inferiori;
  \item $\downarrow$ mobilità articolare.
\end{itemize}

\rubrica{Aspetto psicologico}
\begin{itemize}
  \item $\uparrow$ \textbf{RPE} (\textit{rate of perceived exertion}): la valutazione soggettiva
    che l'atleta esprime subito dopo la seduta (cap.~\ref{cap:profilogiocatore});
  \item $\downarrow$ motivazione;
  \item $\downarrow$ auto-efficacia;
  \item $\downarrow$ ansia.
\end{itemize}
Se persistono, questi segni rimandano non alla fatica di una singola prestazione ma a una
\textbf{stanchezza cronica}, da contrastare con strategie apposite.

\subsection{Che cosa dice la letteratura}

\rubrica{Glicogeno}
\begin{itemize}
  \item riduzione del \textbf{ritmo di lavoro} a fine partita (\textbf{Mohr, Krustrup e Bangsbo},
    2005);
  \item $\downarrow$ \textbf{numero di sprint} (\textbf{Krustrup e coll.}, 2006);
  \item alterazione dell'\textbf{abilità di tiro} (\textbf{Ali e coll.}, 2007).
\end{itemize}

\rubrica{Acidosi intramuscolare}
\begin{itemize}
  \item aumenta dopo le \textbf{fasi intense} di gara, con relazione \textbf{debole ma
    significativa} fra lattato e diminuzione del numero di sprint (\textbf{Krustrup e coll.},
    2003);
  \item $\downarrow$ \textbf{pH muscolare}, che mostra però una relazione \textbf{bassa} con
    l'abbassamento della prestazione (\textbf{Krustrup e coll.}, 2006);
  \item $\uparrow$ RPE (\textbf{Cairns}, 2006).
\end{itemize}

\rubrica{Fosfati intramuscolari e potassio plasmatico}
\begin{itemize}
  \item \textbf{fatica temporanea} (\textbf{Mohr e coll.}, 2005);
  \item $\downarrow$ forza (\textbf{Cairns e Lindinger}, 2008);
  \item $\downarrow$ tempo di esaurimento (\textbf{McKenna e coll.}, 2006).
\end{itemize}

\rubrica{Disidratazione}
\begin{itemize}
  \item $\downarrow$ tempo di esaurimento (\textbf{Jones e coll.}, 2008);
  \item $\uparrow$ RPE (\textbf{Edwards e coll.}, 2007);
  \item $\downarrow$ \textbf{decision making} (\textbf{Cian e coll.}, 2001) --- particolarmente
    rilevante nel calcio a 5, dove il tempo per decidere è brevissimo.
\end{itemize}
È anche il fattore su cui si interviene con più facilità: idratando.

\subsection{Il livello metabolico di partenza}

Studio di \textbf{Ruscello e coll.} (2015), \textit{Acute effect of two different initial heart
rates on testing the repeated sprint ability in young soccer players}: analisi degli effetti acuti
di due diverse percentuali di frequenza cardiaca \textbf{iniziale} sulla prestazione in un test
RSA.

\rubrica{Protocollo}
\begin{itemize}
  \item due riscaldamenti, fino al \textbf{60\,\%} e fino al \textbf{90\,\%} della $FC_{max}$;
  \item 2 serie da 10 ripetizioni di navetta \textbf{15 $+$ 15\,m};
  \item rapporto lavoro:recupero \textbf{1:3}, quello ricavato per la modalità a navetta.
\end{itemize}

\rubrica{Come si raggiungeva la soglia}
In giornate diverse, il soggetto correva uno \textbf{Yo-Yo endurance test} dal primo livello, cioè
da una corsa molto lenta, con un \textbf{cardiofrequenzimetro in telemetria} letto in diretta
dall'operatore. Appena la frequenza toccava la soglia --- 60\,\% o 90\,\% --- il soggetto veniva
\textbf{tolto} dal test e mandato \textbf{subito} a eseguire la prova di RSA.

\begin{table}[htbp]
  \centering
  \small
  \renewcommand{\arraystretch}{1.3}
  \begin{tabularx}{\textwidth}{|X|c|c|X|c|c|}
    \hline
    \textbf{Prova} & \textbf{60\,\%} & \textbf{90\,\%} &
    \textbf{Prova} & \textbf{60\,\%} & \textbf{90\,\%} \\
    \hline
    T1 & 0,00\,\% & 0,00\,\% & T6 & 0,66\,\% & 1,97\,\% \\
    \hline
    T2 & 0,33\,\% & 0,49\,\% & T7 & 1,32\,\% & 2,63\,\% \\
    \hline
    T3 & 0,50\,\% & 1,15\,\% & T8 & 1,49\,\% & 3,78\,\% \\
    \hline
    T4 & 0,83\,\% & 1,31\,\% & T9 & 2,48\,\% & 3,94\,\% \\
    \hline
    T5 & 0,83\,\% & 1,48\,\% & T10 & 2,64\,\% & 4,93\,\% \\
    \hline
  \end{tabularx}
\end{table}

\rubrica{I risultati}
\begin{itemize}
  \item i \textbf{tempi di percorrenza} sono nettamente migliori partendo dal livello metabolico
    più basso, e si allungano partendo dal più alto;
  \item l'\textbf{indice di fatica} partendo dal 90\,\% è circa il \textbf{doppio} a ogni prova, e
    alla decima raggiunge il 4,93\,\% contro il 2,64\,\%. Le rette di regressione danno
    $y = 0,0027x - 0,0039$ ($R^2 = 0,8817$) al 60\,\% e $y = 0,0052x - 0,007$
    ($R^2 = 0,9612$) al 90\,\%;
  \item il \textbf{lattato ematico} è più alto nella condizione al 90\,\% sia prima
    (4,12 contro 2,52\,mmol/l) sia a 3$'$ dalla fine del test (15,02 contro 14,05\,mmol/l).
\end{itemize}

\rubrica{Che cosa se ne ricava}
La RSA che serve in partita \textbf{non} è quella che parte da una frequenza cardiaca di riposo: è
quella che deve restare efficiente \textbf{intorno al 90\,\%} della $FC_{max}$, che è il regime di
gara del calcio a 5. Poiché i \textbf{processi di fatica cambiano} a seconda del livello metabolico
di partenza, sia in allenamento sia in valutazione vanno riprodotte il più fedelmente possibile le
\textbf{condizioni metaboliche della partita}.

\subsection{L'allenamento integrativo di sprint ripetuti}

Studio di \textbf{Soares-Caldeira e coll.} (2014), studio \textbf{randomizzato e controllato}:
aggiungere al normale allenamento precampionato delle sessioni settimanali di sprint ripetuti
migliora la RSA?

\rubrica{Protocollo}
\begin{itemize}
  \item \textbf{13 giocatori} di futsal (22,6 $\pm$ 6,7 anni; 72,8 $\pm$ 8,7\,kg;
    173,2 $\pm$ 6,2\,cm), divisi in due gruppi randomizzati: \textbf{AddT} ($n = 6$), con
    l'allenamento aggiuntivo, e \textbf{NormT} ($n = 7$), con il solo allenamento normale;
  \item valutazione \textbf{pre e post} le 4 settimane di preparazione precampionato;
  \item batteria di test: \textbf{RSA} 6 $\times$ 40\,m a navetta (20 $+$ 20\,m, cambio di
    direzione a 180\textdegree) con 20$''$ di recupero, \textbf{Yo-Yo IR1}, \textbf{SJ} e
    \textbf{CMJ}.
\end{itemize}

\rubrica{I risultati}
\begin{itemize}
  \item in \textbf{entrambi} i gruppi migliorano in modo significativo Yo-Yo IR1, RSA$_{mean}$,
    RSA$_{worst}$ e RSA$_{decrement}$: la capacità mostra dunque un adattamento
    \textbf{rapido};
  \item il guadagno allo Yo-Yo IR1 è del 31,2\,\% nel gruppo con allenamento aggiuntivo contro il
    23,9\,\% dell'altro; sul decremento dell'RSA, però, il gruppo di controllo migliora di più
    ($-$30,4\,\% contro $-$17,8\,\%);
  \item la \textbf{forza esplosiva non migliora}: il CMJ resta invariato in entrambi i gruppi e lo
    squat jump peggiora nel gruppo con l'allenamento aggiuntivo.
\end{itemize}

\rubrica{Che cosa se ne ricava}
L'allenamento integrativo prevedeva sprint ripetuti \textbf{in linea}: allenare la RSA
\textbf{senza} cambi di direzione non basta a migliorare la forza esplosiva. Anche per
l'allenamento, quindi, e non solo per la valutazione, conviene ricorrere alla modalità con
\textbf{cambi di direzione}.
